\subsection{Billion-document estimates}
\label{sec:billion}
The C++ cost models were checked through 8M rows. The bitmap test reaches the full width
of a 1B-row mask, but no complete 1B index was built or searched. The numbers below are
estimates under the same 256-bit score, top 100, one CPU thread and stated distribution of
strong coordinates. A different machine or pattern of memory reads can change the costs.

Start with memory, because it can reject a configuration without estimating its latency.
For the current layout at $N=10^9$ and $d=256$,
\[
M_{\rm codes}=32\ \mathrm{GB},\qquad M_{\rm IDs}=8\ \mathrm{GB}.
\]
A and C need at least 40 GB before directories, containers and query buffers. B retains a
second transposed copy, adding at least 32 GB. Therefore all three fail a 32 GB limit, and B
also fails a 64 GB limit. Changing probes does not remove stored documents. These rejections
apply to the current in-memory layouts; they are not claims about every compressed or disk-based design.

The planning figures below allow room for up to twice the stored row-ID count, a small
allowance for directories and containers, and a configurable 1 GB reserve for running queries.
That reserve is assumed, not measured. Payload counts exclude overhead: if even those bytes
exceed the budget, rejection is certain for this layout. A planned fit is only an estimate. Building an index from full input copies can require more RAM than serving it.

\begin{center}\small
\begin{tabular}{@{}lrrr@{}}\toprule
1B-row path & Payload lower bound & Planned query RAM & 64 GB budget\\\midrule
A, known strong-prefix routing & 40 GB & about 49 GB & Modeled fit\\
C, global 13-bit key & 40 GB & about 49 GB & Modeled fit\\
B, one global cluster & 72 GB & about 83 GB & Rejected\\
B, same prefix route, scan-like leaf & 72 GB & about 81 GB & Rejected\\\bottomrule
\end{tabular}\normalsize\end{center}
All these planning figures fit a 1 TB budget. The model does not claim that a 1 TB machine has
the M5 Max's exact memory behavior. The measured row buffers were much smaller than a 40 GB store.

\subsection{Fixed strong prefix}
With thirteen fixed strong coordinates, the ideal group contains an expected
$10^9/2^{13}\approx122{,}070$ rows. Under the independent-bit assumption, a binomial calculation
describes the chance of having fewer than 100. A Chernoff bound puts that probability below
$10^{-26464}$ in this example. This extremely small bound belongs to the stated distribution;
it is not an embedding-accuracy guarantee.

A can route directly to that group and scan it. C can look up the same key and score its posting
range. Their modeled costs are both roughly 3 ms. The data being read after routing is only a
small part of the corpus, although all codes and IDs are still stored.

B must get the same opportunity: it can use that prefix routing and choose a leaf large enough
to scan the selected group. An all-ones leaf-mask profile shows a similar scoring cost. Thus
B's scan fallback is also estimated at roughly 3 ms, while storing a much larger index. Forcing B to
traverse the entire billion-row bitmap would be an unfair parameter choice in this case.
The model is not precise enough to distinguish a timing gap of a few percent here.
A or C needs less memory, but the estimate does not prove one method is always faster.

\subsection{Changing strong coordinates}
A fixed prefix no longer supplies the same filter. The measured B path can select its columns
from the actual query. Keeping thirteen cuts and scaling the leaf threshold to 160,000 gives
an expected final set of about 122,070 rows. Its mask payload at the last split is
\[
(13+2)\ceil{10^9/64}\times8=1.875\ \mathrm{GB}.
\]
The model fitted to C++ runs estimates about 0.10 s for this global B path. If we double
the part of its cost that grows with corpus size, the estimate becomes about 0.20 s. This
checks the effect of a higher assumed cost; it is not a confidence interval or promised
bound. The model estimates about 25 s for a global scan.

That comparison does not establish that B beats every possible IVF configuration. Let $f$ be
the fraction of the corpus A actually scores, and let $T_{\rm route}$ be the time needed to
obtain those candidates. A's approximate cost is
\[
T_A\approx T_{\rm route}+fN c_{\rm scan}.
\]
Ignoring routing time gives A the most favorable comparison. At the current prices, its
scanned fraction must be below about 0.41\% to undercut the 0.10 s B point, or below about
0.82\% if B's variable cost doubles. It must still meet the recall target at that fraction.
If routing already produces a sufficiently small, accurate pool cheaply, scanning can win.
If it cannot, a global or less heavily clustered B path is a candidate worth testing.

\fig{billion-break-even}{A conditional break-even calculation at 1B rows with sufficient RAM. The scan line excludes routing time, favoring A. B is one fixed global operating point, not a retuned curve. Routing recall and the concentrated-query assumption must be checked separately.}{.94}

B can also use a router. In that case, substitute the physical cluster widths and visited
work into the earlier sum over clusters. The graph does not solve that larger optimization
problem or assume a relevant document is always in the probed clusters.

For C, the number of strong coordinates actually captured by its key matters. Under
Equation~\eqref{eq:captured-strong} and the stated bound conditions, a query capturing none
can still require scoring nearly all rows; capturing one leaves roughly half. A fixed short
key can therefore remain expensive even when it has many occupied entries and few empty
lookups. It is not enough to argue that finding a nonempty key is cheap.

\subsection{What the tests support}
The completed tests establish three different outcomes within their declared spaces:
\begin{itemize}
\item On the real Nomic sample, no alternative speedup was established. Even a one-bit key
control did not improve latency. Recall targets must retain their sampling uncertainty.
\item With a fixed strong prefix, direct routing and integrated key lookup are closely related
and both avoid most scoring. C had a small measured implementation advantage in the controlled runs.
\item With changing strong coordinates, B's query-dependent column order produced a large,
independently checked advantage in the constructed data. Increasing the leaf size with corpus
size preserved a useful path; keeping a tiny fixed leaf would not.
\end{itemize}
At a billion rows, the minimum memory counts are exact for this layout. Query times and
the best routing settings remain estimates under stated assumptions. The formulas give
settings to test and the conditions under which costs would be equal. A finite parameter
grid, or an embedding model name, cannot identify the best setting for every workload.

The current calculation is \path{experiments/project_costs.py}; its inputs and all scenarios
are saved in \path{results/projected-cases-2026-09-11/}. The earlier
\path{reports/search-math/numerical-cases/calculate.py} is preserved as a historical assumed-cost
model. Its extra-posting and hash-slot layout must not be substituted for the measured native C layout.
